% TOP_PROBLEM: {"id": 2625, "title": "Kourovka Notebook Problem 21.116", "category_id": 17, "category": {"id": 17, "name": "group_theory", "display_name": "Group Theory", "description": "Problems about groups, group actions, representations, and related algebraic structures.", "slug": "group-theory", "order_index": 17, "created_at": "2026-07-31T15:26:25.671Z"}, "status": "open", "problem_number": "KOU-21.116", "set_id": 10}
% FIRST_POSTED: 2026-10-01
% ENTRY_KIND: statement_only
% UnsolvedMath ID 2625; source catalogue code KOU-21.116
% !TeX program = lualatex
% Definition and sources only; this file is not an LLM solution attempt.
\documentclass[11pt,a4paper]{article}
\usepackage[margin=25mm]{geometry}
\usepackage{fontspec}
\setmainfont{lmroman10-regular.otf}[BoldFont=lmroman10-bold.otf,
 ItalicFont=lmroman10-italic.otf,BoldItalicFont=lmroman10-bolditalic.otf]
\usepackage{amsmath,amssymb,mathtools}
\usepackage{unicode-math}
\setmathfont{latinmodern-math.otf}
\AtBeginDocument{\let\setminus\smallsetminus}
\usepackage{xurl,hyperref}
\hypersetup{colorlinks=true,urlcolor=blue}
\setcounter{secnumdepth}{0}
\setlength{\parindent}{0pt}
\setlength{\parskip}{5pt}
\setlength{\emergencystretch}{3em}
\begin{document}
\section*{Kourovka Notebook Problem 21.116}\label{2625}
{\small UnsolvedMath ID 2625 · KOU-21.116 · Group Theory · Kourovka Notebook - New Problems, Issue 21}
\subsection{Definitions and mathematical statement}
Let $T$ be an infinite rooted, locally finite, spherically homogeneous tree:
every vertex at level $n$ has the same number $m_n\geq2$ of children.
Suppose $G$ acts faithfully by root-preserving automorphisms, transitively on
each level. The rigid stabilizer $\operatorname{rist}_G(v)$ consists of elements
acting trivially outside the subtree below $v$. Put
$\operatorname{rist}_G(n)=\prod_{v\text{ at level }n}\operatorname{rist}_G(v)$;
the factors commute because their subtrees are disjoint.
The group is branch if it has such an action and
$[G:\operatorname{rist}_G(n)]<\infty$ for every $n$.

Treat $G$ as a discrete group. Let $C_b^k(G;\mathbb R)$ be the vector space
of bounded functions $f:G^{k+1}\to\mathbb R$ invariant under simultaneous
left translation of all arguments. The homogeneous coboundary is
\[
(df)(g_0,\ldots,g_{k+1})=
\sum_{i=0}^{k+1}(-1)^i f(g_0,\ldots,\widehat{g_i},\ldots,g_{k+1}).
\]
Its cohomology $H_b^k(G;\mathbb R)=\ker d/\operatorname{im}d$ is bounded
cohomology with trivial real coefficients; the image is not replaced by its
closure. The group is boundedly acyclic if this vanishes for all $k\geq1$.

Is every branch group boundedly acyclic? In cocycle language, every bounded
invariant cocycle of positive degree would have to be the coboundary of a
bounded invariant cochain. No amenability assumption is imposed.
\subsection{Short English statement}
Does every branch group have vanishing bounded cohomology with trivial real coefficients in all positive degrees?
\subsection{Sources}
\begin{itemize}
\item[S1] ULAM.AI and the UnsolvedMath Contributors, supplied problems.json, record 2625 (KOU-21.116), Kourovka Notebook - New Problems, Issue 21. Catalogue identity and source transcription; CC BY 4.0.
\url{https://huggingface.co/datasets/ulamai/UnsolvedMath}.
\item[S2] The Kourovka Notebook, 21st edition, new problems, Problem 21.116. Expanded formulation of the supplied catalogue statement.
\url{https://arxiv.org/pdf/1401.0300v47}.
\item[S3] E. Schesler, Some problems in geometric group theory, Problem 2.2; discusses the finitely generated special case.
\url{https://eduardschesler.de/Questions___homepage.pdf}.
\end{itemize}
\end{document}
